% Options for packages loaded elsewhere
\PassOptionsToPackage{unicode}{hyperref}
\PassOptionsToPackage{hyphens}{url}
\documentclass[
  ukrainian,
]{article}
\usepackage{xcolor}
\usepackage[margin=2.5cm]{geometry}
\usepackage{amsmath,amssymb}
\setcounter{secnumdepth}{-\maxdimen} % remove section numbering
\usepackage{iftex}
\ifPDFTeX
  \usepackage[T1]{fontenc}
  \usepackage[utf8]{inputenc}
  \usepackage{textcomp} % provide euro and other symbols
\else % if luatex or xetex
  \usepackage{unicode-math} % this also loads fontspec
  \defaultfontfeatures{Scale=MatchLowercase}
  \defaultfontfeatures[\rmfamily]{Ligatures=TeX,Scale=1}
\fi
\usepackage{lmodern}
\ifPDFTeX\else
  % xetex/luatex font selection
\fi
% Use upquote if available, for straight quotes in verbatim environments
\IfFileExists{upquote.sty}{\usepackage{upquote}}{}
\IfFileExists{microtype.sty}{% use microtype if available
  \usepackage[]{microtype}
  \UseMicrotypeSet[protrusion]{basicmath} % disable protrusion for tt fonts
}{}
\makeatletter
\@ifundefined{KOMAClassName}{% if non-KOMA class
  \IfFileExists{parskip.sty}{%
    \usepackage{parskip}
  }{% else
    \setlength{\parindent}{0pt}
    \setlength{\parskip}{6pt plus 2pt minus 1pt}}
}{% if KOMA class
  \KOMAoptions{parskip=half}}
\makeatother
\usepackage{graphicx}
\graphicspath{{/Users/kalinicenkovitalijmikolajovic/.codex/worktrees/5c15/Script_for_translate_book/Output_Final/}}
\makeatletter
\newsavebox\pandoc@box
\newcommand*\pandocbounded[1]{% scales image to fit in text height/width
  \sbox\pandoc@box{#1}%
  \Gscale@div\@tempa{\textheight}{\dimexpr\ht\pandoc@box+\dp\pandoc@box\relax}%
  \Gscale@div\@tempb{\linewidth}{\wd\pandoc@box}%
  \ifdim\@tempb\p@<\@tempa\p@\let\@tempa\@tempb\fi% select the smaller of both
  \ifdim\@tempa\p@<\p@\scalebox{\@tempa}{\usebox\pandoc@box}%
  \else\usebox{\pandoc@box}%
  \fi%
}
% Set default figure placement to htbp
\def\fps@figure{htbp}
\makeatother
\ifLuaTeX
\usepackage[bidi=basic,shorthands=off,provide=*]{babel}
\else
\usepackage[bidi=default,shorthands=off,provide=*]{babel}
\fi
\ifLuaTeX
  \usepackage{selnolig} % disable illegal ligatures
\fi
\setlength{\emergencystretch}{3em} % prevent overfull lines
\providecommand{\tightlist}{%
  \setlength{\itemsep}{0pt}\setlength{\parskip}{0pt}}
\usepackage{bookmark}
\IfFileExists{xurl.sty}{\usepackage{xurl}}{} % add URL line breaks if available
\urlstyle{same}
\hypersetup{
  pdflang={uk},
  hidelinks,
  pdfcreator={LaTeX via pandoc}}

\author{}
\date{}

\begin{document}

\section{Теплопровідність у напівнескінченному
середовищі}\label{ux442ux435ux43fux43bux43eux43fux440ux43eux432ux456ux434ux43dux456ux441ux442ux44c-ux443-ux43dux430ux43fux456ux432ux43dux435ux441ux43aux456ux43dux447ux435ux43dux43dux43eux43cux443-ux441ux435ux440ux435ux434ux43eux432ux438ux449ux456}

Розглянемо глибокий резервуар з рідиною, і нехай бічна поверхня
резервуара буде теплоізольованою. Нехай \texttt{u\_0} буде початковою
температурою рідини, а температура повітря над рідиною ---

\pandocbounded{\includegraphics[keepaspectratio]{images/_page_101_Picture_7.jpeg}}

Рис. 13.2. Схематичне представлення проблеми теплопровідності: a ---
тепло тече всередину, якщо ви(0, t) \textless{} 0, і назовні, якщо ви(0,
t) \textgreater{} 0; b --- резервуар настільки глибокий, що гранична
умова внизу не впливає на розв'язок при значеннях x.

кістка дорівнює нулю (орієнтир температур) (див. рис. 13.2). Наше
завдання --- знайти температурне поле в рідині на різних глибинах і в
різний час. Іншими словами, нам потрібно розв'язати проблему

(13.6)

\begin{verbatim}
\begin{aligned}
(\text{УЧП})\quad & u_t = u_{xx}, && 0 < x < \infty,\; 0 < t < \infty, \\
(\Gamma \text{У})\quad & u_x(0,t) - u(0,t) = 0, && 0 < t < \infty, \\
(\text{НУ})\quad & u(x,0) = u_0, && 0 < x < \infty.
\end{aligned}
\end{verbatim}

Щоб розв'язати цю задачу, ми застосуємо перетворення Лапласа до змінної
t. Зверніть увагу, що можна також застосувати перетворення Лапласа до
змінної x, оскільки воно змінюється від 0 до
\texttt{\textbackslash{}infty} . Після застосування перетворення

Ми отримуємо звичайне диференціальне рівняння за змінною X.

(13.7)

\begin{verbatim}
(ODY) \quad sU(x) - u_0 = \frac{d^2U}{dx^2}, \quad 0 < x < \infty,
(\Gamma Y) \quad \frac{dU}{dx}(0) = U(0)
\end{verbatim}

\begin{verbatim}
(Ми перетворили лише рівняння та граничну умову, а не початкову умову.) Маємо звичайне диференціальне рівняння другого порядку, але з однією граничною умовою при x=0. Насправді існує ще одна гранична умова. Він визначається фізичними міркуваннями і фіксується як U(x) \to 0 , коли x \to +\infty . Зверніть увагу, що для спрощення позначення ми опустили параметр s із аргументів функції U(x) всюди, оскільки рівняння (13.7) є диференціальним рівнянням лише відносно x.
\end{verbatim}

Щоб розв'язати задачу (13.7), запишіть загальне рішення (загальний
розв'язок однорідного + часткового розв'язку гетерогенного) TAC:

\begin{verbatim}
U(x) = c_1 e^{V_{5x}} + c_2 e^{-V_{5x}} + u_0/s.
\end{verbatim}

Підставляючи цей вираз у граничну умову (13.7), можна визначити
константи \texttt{c\_1} та \texttt{c\_2} (одразу зрозуміло, що
\texttt{c\_3\ =\ 0},

\pandocbounded{\includegraphics[keepaspectratio]{images/_page_102_Figure_7.jpeg}}

Рис. 13.3. Граф ймовірнісного інтеграла (eri) та його комплементарної
функції (eric).

\begin{verbatim}
Інакше температура буде зростати нескінченно зі збільшенням координати X). Визначаючи сталу c_2 з граничної умови при x=0, отримуємо фінальний вираз для U\left(x\right) (13.8)
\end{verbatim}

\begin{verbatim}
U(x) = u_0 \left\{ \frac{e^{-V \widehat{s} x}}{s(V \widehat{s} + 1)} \right\} + \frac{1}{s}.
\end{verbatim}

\begin{verbatim}
Залишається зробити останній крок. Щоб визначити температурне поле, потрібно обчислити u\left(x,\,t\right)
\end{verbatim}

\begin{verbatim}
u(x, t) = \mathcal{L}^{-1}[U(x, s)]
\end{verbatim}

(ми повернулися до написання всіх аргументів функції \texttt{U(x,\ s)}).
Щоб повернути перетворення Лапласа, ми використаємо таблиці, наведені в
додатку. У результаті ми отримуємо

(13.9)

\begin{verbatim}
u(x, t) = u_0 - u_0 \left[ \text{eric} \left( \frac{x}{2} \sqrt{t} \right) + \text{eric} \left( \sqrt{t} + \frac{x}{2} \sqrt{t} \right) e^{x+t} \right],
\end{verbatim}

де

\begin{verbatim}
\operatorname{eric}(x) = \frac{2}{\sqrt{\pi}} \int_{x}^{\infty} e^{-\xi^{2}} d\xi
\end{verbatim}

є функцією, доповнюючою до ймовірнісного інтеграла. Графіки цієї функції
показані на рис. 13.3. Аналіз цього рішення P H C

\pandocbounded{\includegraphics[keepaspectratio]{images/_page_103_Figure_7.jpeg}}

13.4. Температура всередині напівнескінченного середовища в різний час.

Потребує дуже малих зусиль. Фіг. (13.4) представлено результати
чисельного аналізу, отриманого на комп'ютері, оснащеному плоттером.

\section{ПРИМІТКИ}\label{ux43fux440ux438ux43cux456ux442ux43aux438}

\begin{itemize}
\tightlist
\item
  \begin{enumerate}
  \def\labelenumi{\arabic{enumi}.}
  \tightlist
  \item
    Перетворення Лапласа також можна застосувати до неоднорідного
    диференціального рівняння з частковими похідними (для застосування
    методу розділення змінних рівняння має бути однорідним), згідно з
    ним, загалом, його можна застосовувати до рівнянь із змінними
    коефіцієнтами (метод розділення змінних також може застосовуватися
    до змінних коефіцієнтів). Таблиця 1. 13.2 вказує сфери застосування
    кожного з двох методів.
  \end{enumerate}
\item
  \begin{enumerate}
  \def\labelenumi{\arabic{enumi}.}
  \setcounter{enumi}{1}
  \tightlist
  \item
    Для розв'язання змішаних і граничних задач також використовуються
    перетворення Ганкеля та Мелліна, які, однак, відрізняються від
    перетворення Лапласа в одному аспекті. Перетворення Лапласа замінює
    похідну операцією множення формули
  \end{enumerate}
\end{itemize}

\begin{verbatim}
\mathscr{L}[y'] = s\mathscr{L}[y] - y(0),
\end{verbatim}

Таблиця 13.2 Порівняння методів розділення перетворення Лапласа та
змінних

\textbar{} Метод \textbar{} \textbar{}

\textbar-------------------------\textbar---------------------------\textbar----------------------------------\textbar{}

\textbar{} Перетворення Лапласа \textbar{} Розділення \textbar{}

Гетерогенні PPP \textbar{} Так \textbar{} Так \textbar{}

Гетерогенні ГІ \textbar{} Так \textbar{} Так \textbar{}

Змінні ймовірності \textbar{} Так \textbar{} de' \textbar{}

Нелінійні рівняння \textbar{} Так \textbar{} Так \textbar{}

а перетворення Ганкеля і Мелліна замінюють диференціальний оператор на
множник. Наприклад, для перетворення Ганкеля, визначеного формулою

\begin{verbatim}
H[y] = \int_{0}^{\infty} r J_{o}(\xi r) y(r) dr,
\end{verbatim}

справедливе співвідношення

\begin{verbatim}
H\left[y''\left(r\right)+\frac{1}{r}\ y'\left(r\right)\right]=-\xi^{2}H\left[y\right].
\end{verbatim}

Таким чином можна розв'язати спеціальне диференціальне рівняння з
змінними коефіцієнтами (рівняння Бесселя). 3. Перетворення Лапласа на
змінній t можна інтерпретувати як проекцію площини xt на осі x, що
призводить до У ц ь о м

\pandocbounded{\includegraphics[keepaspectratio]{images/_page_104_Figure_9.jpeg}}

у випадку початкові рівняння, граничні та початкові умови перетворюються
на нове диференціальне рівняння та нові граничні умови (див. діаграму
вище).

\section{ЗАВДАННЯ}\label{ux437ux430ux432ux434ux430ux43dux43dux44f}

\begin{enumerate}
\def\labelenumi{\arabic{enumi}.}
\tightlist
\item
  Переконайтеся, що наступна формула для перетворення часткової похідної
  \texttt{u\_t} є істинною:
\end{enumerate}

\begin{verbatim}
\mathscr{L}\left[u_{t}(x, t)\right] = sU(x, s) - u(x, 0).
\end{verbatim}

\begin{enumerate}
\def\labelenumi{\arabic{enumi}.}
\setcounter{enumi}{1}
\tightlist
\item
  Використати перетворення Лаціаса для розв'язання задачі Коші
\end{enumerate}

\begin{verbatim}
\begin{array}{llll} (\text{УЧП}) & u_t = \alpha^2 u_{xx}, & -\infty < x < \infty, & 0 < t < \infty, \\ (\text{HV}) & u(x, 0) = \sin x, & -\infty < x < \infty. \end{array}
\end{verbatim}

\begin{enumerate}
\def\labelenumi{\arabic{enumi}.}
\setcounter{enumi}{2}
\tightlist
\item
  С помощью преобразования Лапласа по переменной t решите вадачу
\end{enumerate}

\begin{verbatim}
\begin{array}{llll} (\text{УЧ\Pi}) & u_t = u_{x\,\text{v}}, & 0 < x < \infty, & 0 < t < \infty, \\ (\text{ГУ}) & u\left(0,\ t\right) = \sin t, & 0 < t < \infty, \\ (\text{HY}) & u\left(x,\ 0\right) = 0, & 0 \leqslant x < \infty. \end{array}
\end{verbatim}

Дайте фізичну інтерпретацію цієї зали.

\begin{enumerate}
\def\labelenumi{\arabic{enumi}.}
\setcounter{enumi}{3}
\tightlist
\item
  Розв'язати задачу з граничним значенням
\end{enumerate}

(ОДА)

\begin{verbatim}
\frac{d^2U}{dx^2} - sU = A, \quad 0 < x < 1,
(\Gamma Y) \quad
\begin{cases} \frac{dU}{dx} (0) = 0 \\ U(1) = 0. \end{cases}
\end{verbatim}

Лекція 14

\section{ПРИНЦИП
ДЮАМЕЛЯ}\label{ux43fux440ux438ux43dux446ux438ux43f-ux434ux44eux430ux43cux435ux43bux44f}

МЕТА ЛЕКЦІЇ: Показати, як перетворення Лапласа можна використати для
виявлення цікавих особливостей розв'язання диференціальних рівнянь.
Зокрема, за допомогою алгебранічних перетворень образів Лапласа деякого
розв'язку рівняння часткових похідних можливо отримати відомий принцип
Дюамеля. Інтерпретація цього принципу в теорії звичайних диференціальних
рівнянь відома, але ми покажемо, як його використовувати для розв'язання
деяких змішаних задач для UPP.

технології.

Окрім того, що перетворення Лапласа є потужним методом розв'язання
часткових диференціальних рівнянь, воно дозволяє виявити фізичну
сутність отриманих розв'язків. За допомогою трансформації Лапласа ми
проілюструємо дуже важливий і цікавий принцип Дюамеля. Перш ніж перейти
до принципу Дюамеля, розглянемо проблему, яка дуже часто виникає у

\section{Теплопровідність у стрижні, кінці якого підтримуються при
заданій
температурі}\label{ux442ux435ux43fux43bux43eux43fux440ux43eux432ux456ux434ux43dux456ux441ux442ux44c-ux443-ux441ux442ux440ux438ux436ux43dux456-ux43aux456ux43dux446ux456-ux44fux43aux43eux433ux43e-ux43fux456ux434ux442ux440ux438ux43cux443ux44eux442ux44cux441ux44f-ux43fux440ux438-ux437ux430ux434ux430ux43dux456ux439-ux442ux435ux43cux43fux435ux440ux430ux442ux443ux440ux456}

Дуже часто необхідно знайти температуру всередині області, якщо відомо,
що на межі області вона змінюється з часом відповідно до певного закону
(граничні умови, залежні від часу).

На лівому кінці підтримується умова \texttt{u(0,t)}=0, а на правому
кінці температура задається функцією \texttt{f(t)}.

Рис. 14.1. Граничні умови, що залежать від часу.

Наприклад, розглянемо термоізольований стержень, правий кінець якого
підтримується у \texttt{f(t)} (див. рисунок 14.1). Потрібно вирішити цю
проблему

(14.1)

\begin{verbatim}
u_{t} = u_{xx}, \quad 0 < x < 1, \quad 0 < t < \infty,
(\Gamma Y) \quad \begin{cases} u(0, t) = 0, \\ u(1, t) = f(t), \end{cases} \quad 0 < t < \infty,
(HY) \quad u(x, 0) = 0, \quad 0 \le x \le 1.
\end{verbatim}

Сподіваємося, що розв'язання цієї задачі (14.1) буде легким, оскільки ми
знаємо розв'язок простішої задачі, де температура в кінці є сталою:

Насправді, якщо обидві задачі ((14.1) і (14.2)) розв'язуються одночасно
за допомогою перетворення Лапласа, то отримаємо визначний результат:
розв'язок задачі (14.1) виражається через розв'язання задачі (14.2).

Отже, розв'язуючи обидві задачі одночасно, ми отримуємо

\section{Легке завдання (14.2) (Постійні
солдати)}\label{ux43bux435ux433ux43aux435-ux437ux430ux432ux434ux430ux43dux43dux44f-14.2-ux43fux43eux441ux442ux456ux439ux43dux456-ux441ux43eux43bux434ux430ux442ux438}

Застосування перетворення -Лапласа до задачі (14.2)

\begin{verbatim}
\frac{d^2W}{dx^2} - sW(x) = 0,
W(0) = 0,
W(1) = 1/s
\end{verbatim}

.

Пемеріль ODU

\begin{verbatim}
W(x, s) = \frac{1}{s} \left[ \frac{\sinh(x \sqrt{s})}{\sinh(\sqrt{s})} \right].
\end{verbatim}

Пошук оберненого перетворення

\begin{verbatim}
w(x, t) =
= x + \frac{2}{\pi} \sum_{n=1}^{\infty} \frac{(-1)^n}{n} e^{-(n\pi)^2 t} \times
\end{verbatim}

\texttt{\textquotesingle{}\textbackslash{}times\ \textbackslash{}sin(n\textbackslash{}pi\ x)\textquotesingle{}\ (розв\textquotesingle{}язання\ задачі\ з\ постійним\ ГІ).}

\section{Важке завдання (14.1) (Солдати, що залежать від
часу)}\label{ux432ux430ux436ux43aux435-ux437ux430ux432ux434ux430ux43dux43dux44f-14.1-ux441ux43eux43bux434ux430ux442ux438-ux449ux43e-ux437ux430ux43bux435ux436ux430ux442ux44c-ux432ux456ux434-ux447ux430ux441ux443}

Застосування перетворення Лапласа до задачі (14.1)

\begin{verbatim}
\frac{d^2U}{dx^2} - sU(x) = 0,
U\left( \mathbf{0}\right) =0,
U(1) = F(s)
\end{verbatim}

.

\begin{verbatim}
U(x, s) = F(s) \left[ \frac{\sinh(x \sqrt{s})}{\sinh(\sqrt{s})} \right].
\end{verbatim}

Розділіть і помножте на z

\begin{verbatim}
U(x, s) = F(s) \left\{ s \left[ \frac{\sinh(x \sqrt{s})}{s \sinh(s)} \right] \right\}.
\end{verbatim}

Використовуючи співвідношення
\texttt{\textbackslash{}mathscr\{L\}{[}w\_t{]}\ =\ sW\ -\ w(x,\ 0),}
отримуємо

\begin{verbatim}
U(x, s) = F(s) \mathcal{L}[w_t].
\end{verbatim}

Відповідно,

\begin{verbatim}
u(x, t) = \mathcal{L}^{-1} \{ F(s) \mathcal{L}[w_t] \} =
= \mathcal{L}^{-1} [ F(s) ] * \mathcal{L}^{-1}[w_t] =
= f(t) * w_t(t) =
= \int_0^t w_t(x, t-\tau) f(\tau) d\tau =
\end{verbatim}

(після інтегрування частинами)

\begin{verbatim}
= \int_{0}^{t} w(x, t-\tau) f'(\tau) d\tau + f(0) w(x, t).
\end{verbatim}

(Розв'язок задачі з часово-залежними ГІ виражається через розв'язок
задачі з постійними ГІ

розв'язано у задачах з константою

Отже, ми висловили розв'язок задач із граничними умовами, що змінюються
в часі, через розв'язання задачі з постійними граничними умовами.
Відповідні формули такі:

(14.3)

плутанина у наступних формулах. Записати диференціальне рівняння в
похідних у нових координатах (~, т), мы воспользуемся следующими
формулами перехода;

Наприклад, якщо dy/dx = B/2A = 3, то y - 3x = c, а отже
\texttt{\textbackslash{}xi(x,\ y)\ =\ y\ -\ 3x} . Ця функція задовольняє
умову \texttt{{[}\textbackslash{}xi\_x/\textbackslash{}xi\_y{]}\ =\ -3}
і, відповідно, повертає коефіцієнт
\texttt{\textbackslash{}overline\{A\}} до нуля.

Отже, половина роботи виконана: ми знайшли нову координатну
\texttt{\textbackslash{}xi\ =\ \textbackslash{}xi(x,\ y)} , яка
забезпечила рівність \texttt{\textbackslash{}overline\{A\}\ =\ 0} .
Залишається знайти \texttt{\textbackslash{}eta(x,\ y)} таким чином, щоб
перетворити коефіцієнт \texttt{\textbackslash{}overline\{B\}} на нуль.

Час зупинитися тут. Виявляється, що якщо вибрати \texttt{B\^{}2-4AC=0} і
\texttt{\textbackslash{}xi} так, щоб
\texttt{\textbackslash{}overline\{A\}=0} , то коефіцієнт
\texttt{\textbackslash{}overline\{B\}} автоматично стане нулем 1). Тепер
ми побачимо це.

Коефіцієнт \texttt{\textbackslash{}overline\{B\}} задається виразом

\begin{verbatim}
\overline{B}^2 - 4\overline{A}\overline{C} = (B^2 - 4AC) (\xi_x \eta_y - \xi_y \eta_x)^2 - \Pi \rho u u. \ \rho e \partial.
\overline{B} = 2 \left( V \overline{A} \xi_x + V \overline{C} \xi_y \right) \left( V \overline{A} \eta_x + V \overline{C} \eta_y \right).
\end{verbatim}

Функциональная диаграмма на рис. 15.5 делает эти формулы совсем
очевидными. Эта диаграмма особенно полезна при вычислении частных
производных функций и и и по переменным Р И С

\pandocbounded{\includegraphics[keepaspectratio]{images/_page_115_Picture_14.jpeg}}

15.5 Диаграмма, иллюстрирующая функциональную зависимость переменных

\begin{verbatim}
\xi   и   \tau  , поскольку   \xi   зависит и от x, и от t. Переменная   \tau   зависит только от t. Это все, что касается преобразований. Теперь подставим найденные выражения   u_t  ,   u_x   и   u_{vx}   в уравнение (15 2) и получим
\end{verbatim}

\begin{verbatim}
\overline{B} = 2 \sqrt{\overline{A}} \left[ \sqrt{\overline{A}} \, \eta_x + \sqrt{\overline{C}} \, \eta_y \right].
\end{verbatim}

откуда

\begin{verbatim}
u_{xy} + 2u_{xy} + u_{yy} = 0.
\end{verbatim}

.

\begin{verbatim}
Следовательно, новая задача с начальным условием в персменных   \xi   и   \tau   имеет вид
\end{verbatim}

(УЧП)

\begin{verbatim}
\frac{dy}{dx} = -[\xi_x/\xi_y] = B/2A = 1.
\end{verbatim}

\begin{verbatim}
(НУ)   u(\xi, 0) = 1 - H(\xi), \quad -\infty < \xi < \infty.
\end{verbatim}

\begin{verbatim}
Так как   \xi = x   при t = 0, то новые граничные условия совпадают со старыми. Эта задача была решена в лекции 12 с помощью интегрального преобразования Фурье, и ее решение записывается в виде
\end{verbatim}

\begin{verbatim}
y = x + c
\end{verbatim}

\begin{verbatim}
где   \phi(\beta)   — начальное условие. Следовательно, в нашем случае
\end{verbatim}

\begin{verbatim}
\eta = y
\end{verbatim}

После подстановки

\begin{verbatim}
\begin{cases}
 \xi = y - x, \\
 \eta = y
 \end{cases}
\end{verbatim}

получаем окончательный результат

(15.3)

\begin{verbatim}
\overline{C} = A \eta_x^2 + B \eta_x \eta_y + C \eta_y^2 = 1,
\overline{E} = A \eta_{xx} + B \eta_{xy} + C \eta_{yy} + D \eta_x + E \eta_y = 0,
\begin{cases} \frac{1}{2} \left[ 1 + \operatorname{erf} \left( \frac{-\xi}{2\sqrt{D\tau}} \right) \right], & \xi < 0, \\ \frac{1}{2} \operatorname{erfc} \left( \frac{\xi}{2\sqrt{D\tau}} \right), & 0 \leq \xi. \end{cases}
\overline{G} = G = 0.
\overline{A}u_{\xi\xi} + \overline{B}u_{\xi\eta} + \overline{C}u_{\eta\eta} + \overline{D}u_{\xi} + \overline{E}u_{\eta} + \overline{F}u = \overline{G}
\end{verbatim}

Графики этого решения для различных моментов времени изображены на рис.
15.6. Осталось сделать последний шаг --- записать решение нашей задачи в
координатах x и t. Подставляя

\begin{verbatim}
u_{\eta} = f(\xi),
,   \tau = t
\end{verbatim}

в формулы (15.3), получаем

(15.4)

\begin{verbatim}
u(x, t) = \begin{cases} \frac{1}{2} \left[ 1 + \operatorname{erf}\left(\frac{Vt - x}{2\sqrt{Dt}}\right) \right], & Vt > x, \\ \frac{1}{2} \operatorname{erfc}\left(\frac{x - Vt}{2\sqrt{Dt}}\right), & Vt \leq x. \end{cases}
\end{verbatim}

Перед нами розв'язок задачі конвективної дифузії (15.2). Профіль
концентрації рухається праворуч зі швидкістю \texttt{V} і водночас
розмивається дифузією з коефіцієнтом \texttt{D}.

\paragraph{ЗАВДАННЯ}\label{ux437ux430ux432ux434ux430ux43dux43dux44f-1}

\begin{enumerate}
\def\labelenumi{\arabic{enumi}.}
\tightlist
\item
  Розв'яжіть задачу Коші
\end{enumerate}

\begin{verbatim}
\begin{array}{ll} u_t = u_{xx} - 2u_x, & -\infty < x < \infty, \ 0 < t < \infty, \\ u(x, 0) = \sin x, & -\infty < x < \infty. \end{array}
\end{verbatim}

\begin{enumerate}
\def\labelenumi{\arabic{enumi}.}
\setcounter{enumi}{1}
\tightlist
\item
  Знайдіть розв'язок задачі Коші для рівняння конвективної дифузії
\end{enumerate}

\begin{verbatim}
\begin{array}{ll} u_t = u_{xx} - 2u_x, & -\infty < x < \infty, \ 0 < t < \infty, \\ u(x, 0) = e^x \sin x, & -\infty < x < \infty. \end{array}
\end{verbatim}

\begin{enumerate}
\def\labelenumi{\arabic{enumi}.}
\setcounter{enumi}{2}
\tightlist
\item
  Знайдіть розв'язок задачі переносу, використовуючи перетворення
  координат з лекції 8.
\end{enumerate}

\subsection{Приведення еліптичних рівнянь до канонічної
форми}\label{ux43fux440ux438ux432ux435ux434ux435ux43dux43dux44f-ux435ux43bux456ux43fux442ux438ux447ux43dux438ux445-ux440ux456ux432ux43dux44fux43dux44c-ux434ux43e-ux43aux430ux43dux43eux43dux456ux447ux43dux43eux457-ux444ux43eux440ux43cux438}

Розглянемо загальне рівняння

\begin{verbatim}
Au_{xx} + Bu_{xy} + Cu_{yy} + Du_x + Eu_y + Fu = G
\end{verbatim}

але тепер у випадку \texttt{B\^{}2\ -\ 4AC\ \textless{}\ 0}. Переходом
до нових незалежних змінних ми хочемо звести його до форми

\begin{verbatim}
u_{\xi\xi} + u_{\eta\eta} = \varphi(\xi, \eta, u, u_{\xi}, u_{\eta}).
\end{verbatim}

Щоб знайти ці нові змінні, виконуємо ті самі обчислення, що й у двох
попередніх випадках, отримуємо перетворене рівняння

\begin{verbatim}
\overline{A}u_{\xi\xi} + \overline{B}u_{\xi\eta} + \overline{C}u_{\eta\eta} + \overline{D}u_{\xi} + \overline{E}u_{\eta} + \overline{F}u = \overline{G},
\end{verbatim}

і вимагаємо, щоб \texttt{\textbackslash{}xi} та
\texttt{\textbackslash{}eta} задовольняли умовам
\texttt{\textbackslash{}overline\{A\}\ =\ \textbackslash{}overline\{C\}}
і \texttt{\textbackslash{}overline\{B\}\ =\ 0}. Безпосередньо знайти
\texttt{\textbackslash{}xi} та \texttt{\textbackslash{}eta} тут
складніше, тому шукаємо перетворення як композицію двох кроків.

\section{Перетворення
1}\label{ux43fux435ux440ux435ux442ux432ux43eux440ux435ux43dux43dux44f-1}

Спочатку введемо комплексні координати \texttt{\textbackslash{}xi} та
\texttt{\textbackslash{}eta}, щоб надати рівнянню вигляду

\begin{verbatim}
u_{\xi\eta} = \psi(\xi, \eta, u, u_{\xi}, u_{\eta}).
\end{verbatim}

Для цього розв'язуємо характеристичні рівняння

\begin{verbatim}
\frac{dy}{dx} = \frac{B - \sqrt{B^2 - 4AC}}{2A},
\end{verbatim}

де \texttt{B\^{}2\ -\ 4AC\ \textless{}\ 0}, і отримуємо

\begin{verbatim}
\xi(x, y) = \text{const}, \qquad \eta(x, y) = \text{const}.
\end{verbatim}

\section{Перетворення
2}\label{ux43fux435ux440ux435ux442ux432ux43eux440ux435ux43dux43dux44f-2}

Далі переходимо від \texttt{(\textbackslash{}xi,\textbackslash{}eta)} до
\texttt{(\textbackslash{}alpha,\textbackslash{}beta)} за формулами

\begin{verbatim}
\alpha = \frac{\xi + \eta}{2}, \qquad \beta = \frac{\xi - \eta}{2i}.
\end{verbatim}

У результаті отримуємо канонічну форму еліптичного рівняння.

\section{Зведення рівняння y\^{}2u\_\{xx\} + x\^{}2u\_\{yy\} = 0 до його
канонічної
форми}\label{ux437ux432ux435ux434ux435ux43dux43dux44f-ux440ux456ux432ux43dux44fux43dux43dux44f-y2u_xx-x2u_yy-0-ux434ux43e-ux439ux43eux433ux43e-ux43aux430ux43dux43eux43dux456ux447ux43dux43eux457-ux444ux43eux440ux43cux438}

Розглянемо рівняння

у яких \texttt{A=y\^{}2} , B=0, \texttt{C=x\^{}2} , D=E=F=G=0.
Дискримінантний \texttt{B\^{}2-4AC} дорівнює \texttt{-4x\^{}2y\^{}2} , і
ми зведемо це рівняння до канонічної форми в першому квадранті
x\textgreater0, y\textgreater0.

КРОК 1 (Перше перетворення)

Запишемо характеристичні рівняння

\begin{verbatim}
\frac{dy}{dx} = \frac{B + \sqrt{B^2 - 4AC}}{2A} = \frac{\sqrt{-4x^2y^2}}{2y^2} = i\frac{x}{y}
\end{verbatim}

сделайте замену переменных \texttt{\textbackslash{}xi\ =\ x\ -\ Vt} ,
\texttt{\textbackslash{}tau\ =\ t} . В нашем случае можно в явном виде
вычислить интеграл

\begin{verbatim}
y^2 + ix^2 = \text{const},
\end{verbatim}

Такой интеграл уже встречался в лекции 12. Может быть, читателю будет
удобнее обратиться к таблицам преобразования Фурье.

\section{ОДНОМЕРНОЕ ВОЛНОВОЕ УРАВНЕНИЕ (ГИПЕРБОЛИЧЕСКИЕ
УРАВНЕНИЯ)}\label{ux43eux434ux43dux43eux43cux435ux440ux43dux43eux435-ux432ux43eux43bux43dux43eux432ux43eux435-ux443ux440ux430ux432ux43dux435ux43dux438ux435-ux433ux438ux43fux435ux440ux431ux43eux43bux438ux447ux435ux441ux43aux438ux435-ux443ux440ux430ux432ux43dux435ux43dux438ux44f}

ЦЕЛЬ ЛЕКЦИИ: Показать, что колебания сгруны описываются уравнением

Від

и что это уравнение --- следствие законов Ньютона. Рассматриваются и
некоторые другие виды волнового уравнения:

\begin{verbatim}
\eta(x, y) = y^2 - ix^2. (Неважливо, яка з двох функцій називається \xi , а яка — \eta ; ви можете поміняти їх місцями, якщо хочете.) Після цього перетворення початкове рівняння набуває вигляду
\end{verbatim}

Ми не звертатимемо уваги на це рівняння (яке є комплексним гіперболічним
рівнянням), а одразу перейдемо до другого перетворення.

КРОК 2. (Друга трансформація)

При здійсненні другої трансформації ми маємо

\begin{verbatim}
\begin{split} u_{tt} &= \alpha^2 u_{xx} + F\left(x, \ t\right), \\ u_{tt} &= \alpha^2 u_{xx} - \beta u_t - \lambda u + F\left(x, \ t\right). \end{split}
\end{verbatim}

\texttt{(реальна\ частина\ \textquotesingle{}\textbackslash{}xi\textquotesingle{}\ і\ \textquotesingle{}\textbackslash{}eta\textquotesingle{}\ ),\ \textquotesingle{}\textbackslash{}beta\ =\ \textbackslash{}frac\{\textbackslash{}xi\ -\ \textbackslash{}eta\}\{2\textbackslash{}iota\}\ =\ x\^{}2\textquotesingle{}\ (уявна\ частина\ \textquotesingle{}\textbackslash{}xi\textquotesingle{}\ і\ \textquotesingle{}\textbackslash{}eta\textquotesingle{}\ ).}

Щоб спростити нотацію, найкраще перепозначити
\texttt{(\textbackslash{}alpha,\textbackslash{}beta)}
\texttt{(\textbackslash{}xi,\textbackslash{}eta)} окремо і записати
отримане \texttt{\textbackslash{}epsilon} перетворення як

КРОК 3 (Отримання нового рівняння)

\begin{verbatim}
Нову канонічну форму можна отримати, обчисливши всі коефіцієнти \overline{A}, \overline{B}, \overline{C}, \overline{D}, \overline{E}, \overline{F} і \overline{G} у рівнянні
\end{verbatim}

До сих пор мы занкмались физическими процессами, которые описываются
одномерными параболическими уравнениями (диффузионные задачи). Теперь мы
приступим к изучению еще одного класса уравнений с частными
производными---гиперболических уравнений. Начнем с одномерного волнового
уравнения, которое описывает, в частности, поперечные колебания струны.

\section{Колебания
струны}\label{ux43aux43eux43bux435ux431ux430ux43dux438ux44f-ux441ux442ux440ux443ux43dux44b}

Рассмотрим малые колебания струны с закрепленными концами. Предположим,
что струна туго натянута, сделана из однородного материала и колеблется
в одной плоскости (рис. 16.1).

\pandocbounded{\includegraphics[keepaspectratio]{images/_page_119_Figure_10.jpeg}}

РИС. 16.1. Поперечные колебания струны: a --- смещение струны от
положения равновесия.

Действием сил тяжести на струну будем пренебрегать. Для построения
математической модели колебаний струны рассмотрим все силы, действующие
на небольшой участок струны (рис. 16.2). Оказывается, что волновое
уравнение---это не что иное, как закон движения Ньютона (изменение
импульса равно сумме действующих сил), примененный к струне. Глядя на
рис. 16.1, мы можем представить те силы, которые действуют на струну в
направлений, перпендикулярном оси х.

\pandocbounded{\includegraphics[keepaspectratio]{images/_page_120_Figure_2.jpeg}}

\texttt{РИС.\ 16.2.\ Малый\ участок\ \ \textquotesingle{}(x,\ x\ +\ \textbackslash{}Delta\ x)\textquotesingle{}\ \ колеблющейся\ струны.}

\begin{enumerate}
\def\labelenumi{\arabic{enumi}.}
\tightlist
\item
  Сила, обусловленная натяжением струны (
  \texttt{\textbackslash{}alpha\^{}2\ u\_\{xx\}} ). Величина поперечной
  компоненты силы натяжения находится по формуле
\end{enumerate}

Сила =

\begin{verbatim}
за формулами (41.5) з \xi = y^2 та \eta = x^2 . Зробивши це, отримуємо канонічну форму еліптичного рівняння
\end{verbatim}

\texttt{\textquotesingle{}\textbackslash{}simeq\ T\ {[}u\_x(x\ +\ \textbackslash{}Delta\ x,\ t)\ -\ u\_x(x,\ t){]}.\textquotesingle{}\ 2.\ Внешняя\ сила\ \textquotesingle{}F\ (x,\ t)\textquotesingle{}.}

Внешняя сила \texttt{F(x,\ t)}, приложенная к струне, может произвольным
образом зависеть от x и t. Некоторые примеры внешних сил приводятся
ниже:

а) гравитационная сила \texttt{F(x,\ t)} = -mg;

б) импульсы, распределенные вдоль струны и действующие на нее в
различные моменты времени:

в) силы, возникающие при действии звуковой волны на кожу барабана (позже
мы подробно изучим колебания мембраны, которые описываются волновым
уравнением).

\begin{enumerate}
\def\labelenumi{\arabic{enumi}.}
\setcounter{enumi}{2}
\item
  Силы трения, действующие на струну (---
  \texttt{\textbackslash{}beta\ u\_t} ). Если струна колеблется в среде,
  то возникает сила сопротивления, которая пропорциональна скорости
  \texttt{u\_t} .
\item
  Возвращающая сила (- үи).
\end{enumerate}

Эта сила направлена противоположно смещению струны. Если смещение u
положительно, то сила отрицательна.

Если применить теперь уравнения движения Ньютона к небольшому участку
струны, то получим

\section{ПРИМІТКИ}\label{ux43fux440ux438ux43cux456ux442ux43aux438-1}

\begin{enumerate}
\def\labelenumi{\arabic{enumi}.}
\tightlist
\item
  Аналогічно можна класифікувати лінійні рівняння другого порядку з
  \emph{трьох або більше змінних}, але для цього потрібен матричний
  аналіз. Наприклад, рівняння другого порядку з трьома змінними
\end{enumerate}

\texttt{где\ \ \textquotesingle{}\textbackslash{}rho\textquotesingle{}\ \ —\ линейная\ плотность\ струны.\ Разделим\ обе\ части\ этого\ сротношения\ на\ \ \textquotesingle{}\textbackslash{}Delta\ x\textquotesingle{}\ \ и\ устремим\ \ \textquotesingle{}\textbackslash{}Delta\ x\textquotesingle{}\ \ —\ 0.\ В\ результате\ получим\ хорошо\ известное\ телеграфное\ уравнение}

(16.1)

згідно з цією класифікацією, буде параболічним, а рівняння

В этом уравнении следовало бы разделить \texttt{\textbackslash{}alpha} ,
\texttt{\textbackslash{}beta} , \texttt{\textbackslash{}gamma} и
\texttt{F(x,t)} на \texttt{\textbackslash{}rho} , но мы ради простоты
сохраним для этих величин прежние обозначения. Искомое уравнение
получено. Приведем интерпретацию простейшего волнового уравнения

\begin{itemize}
\tightlist
\item
  Гіперболічно.
\end{itemize}

\begin{enumerate}
\def\labelenumi{\arabic{enumi}.}
\setcounter{enumi}{1}
\tightlist
\item
  Наш інтерес до класифікації диференціальних рівнянь у похідних
  пов'язаний із тим, що три основні типи рівнянь описують три різні типи
  фізичних явищ, і ми хотіли б отримати математичну класифікацію цих
  трьох класів фізичних явищ.
\end{enumerate}

\section{ЗАВДАННЯ}\label{ux437ux430ux432ux434ux430ux43dux43dux44f-2}

\begin{enumerate}
\def\labelenumi{\arabic{enumi}.}
\tightlist
\item
  Які з цих параболічних і еліптичних рівнянь записані у канонічній
  формі?
\end{enumerate}

\begin{verbatim}
a) u_t = u_{xx} + hu,
б) u_{xy} + u_{xx} + 3u = \sin x,
в) u_{xx} + 2u_{yy} = 0,
г) u_{xx} = \sin^2 x?
\end{verbatim}

\begin{enumerate}
\def\labelenumi{\arabic{enumi}.}
\setcounter{enumi}{1}
\item
  Перетворіть параболічне рівняння \texttt{u\_\{xx\}} + 2u\_\{xy\} +
  \texttt{u\_\{yy\}} + u = 2 до канонічної форми.
\item
  Перетворіть еліптичне рівняння \texttt{u\_\{xx\}} + 2u\_\{yy\} +
  x\^{}2 \texttt{u\_x\ =\ e\^{}\{-x\^{}2/2\}}до канонічної форми.
\end{enumerate}

\section{МЕТОД МОНТЕ-КАРЛО
(ВСТУП)}\label{ux43cux435ux442ux43eux434-ux43cux43eux43dux442ux435-ux43aux430ux440ux43bux43e-ux432ux441ux442ux443ux43f}

МЕТА ЛЕКЦІЇ: Пояснити основні ідени, що лежать в основі методу
Монте-Карло, і показати, як цей метод можна використовувати для
розв'язання різних задач. Виявляється, що за допомогою азартних ігор
можна змоделювати (загалом, на комп'ютері) наближені розв'язки реальних
задач. Як простий приклад можна вказати обчислення інтеграла

основанную на интуитивных соображениях.

\section{Интуитивная интерпретация волнового
уравнения}\label{ux438ux43dux442ux443ux438ux442ux438ux432ux43dux430ux44f-ux438ux43dux442ux435ux440ux43fux440ux435ux442ux430ux446ux438ux44f-ux432ux43eux43bux43dux43eux432ux43eux433ux43e-ux443ux440ux430ux432ux43dux435ux43dux438ux44f}

У читателя может возникнуть вопрос: почему уравнение вида (16.2) должно
описывать нечто такое, что похоже на колебания скрипичной струны? Чтобы
разобраться в этом, необходимо понять,

\pandocbounded{\includegraphics[keepaspectratio]{images/_page_121_Picture_9.jpeg}}

РИС. 16 3. Интерпретация волнового уравнения
\texttt{u\_\{tt\}\ =\ \textbackslash{}alpha\^{}2\ u\_\{xx\}}

что величина \texttt{u\_\{tt\}} является вертикальным ускорением струны
в точке x. Следовательно, уравнение (16.2) можно интерпретировать
следующим образом: ускорение струны, обусловленное натяжением, в каждой
точке тем больше, чем больше вогнутость струны \texttt{u\_\{xx\}} в
данной точке (коэффициентом пропорциональности служит константа
\texttt{\textbackslash{}alpha\^{}2\ =\ T/\textbackslash{}rho} ) (см.
рис. 133).

\section{ЗАМЕЧАНИЯ}\label{ux437ux430ux43cux435ux447ux430ux43dux438ux44f}

\texttt{1.\ Волновое\ уравнение\ \ \textquotesingle{}u\_\{tt\}\ =\ \textbackslash{}alpha\^{}2\ u\_\{xx\}\textquotesingle{}\ \ описывает\ также\ продольные\ и\ крутильные\ колебания\ стержня.\ В\ продольных\ колебаниях\ смещения\ параллельны\ стержню,\ а\ через\ \ \textquotesingle{}u\textbackslash{}left(x,\textbackslash{},t\textbackslash{}right)\textquotesingle{}\ \ обозначается\ величина\ продольного\ смещения\ относительно\ равновесного\ положения.\ Такие\ колебания\ возникают,\ например,\ при\ ударе\ молотком\ по\ торцу\ стержня.\ Аналогично\ описываются\ крутиль-}

ные колебания

\$\$

кидаючи дротик у одиничний квадрат
\texttt{\textbackslash{}\{(x,\ y):\ 0\ \textless{}\ x\ \textless{}\ 1,\ 0\ \textless{}\ y\ \textless{}\ 1\textbackslash{}\}}
Після приблизно 100 кидків відносна частина дротиків, що опускається
нижче кривої \texttt{y\ =\ x\^{}2} , вважається приблизним значенням
інтеграла. Узагальнення такої гри (з випадковим киданням дрогіка) завжди
містить процедуру генерації послідовності випадкових чисел.

У цій лекції ми коротко опишемо метод Монте-Карло, а в наступній лекції
покажемо, як його використовувати для розв'язання диференціальних
рівнянь у похідних

Методи Монте-Карло зазвичай називають групою методів розв'язання
детермінованих задач (тобто задач без випадковості), у яких суттєво
використовуються елементи випадковості. Загальна філософія етнічних
методів показана на рис. 42 1.

Ймовірнісна гра

Деторміноване завдання \texttt{\textbackslash{}rho} --- розв'язання
задачі (наприклад, обчислення інтеграла, розв'язання UBP тощо).

Результат = \texttt{\textbackslash{}hat\{p\}} Наближення p Рішення = p
Рис. 42.1. Загальна теорія методу Монте-Карло.

\end{document}
